\chapter{Varicella}
\index{Varicella}

\section*{Definition and Overview}
Varicella, commonly known as chickenpox, is a highly contagious infectious disease caused by primary infection with the varicella-zoster virus (VZV), a double-stranded DNA virus belonging to the Herpesviridae family (specifically human alphaherpesvirus 3). The disease is clinically characterized by a pruritic vesicular skin rash, fever, and systemic malaise. Historically considered a universal rite of passage during childhood, varicella represents a significant public health focus due to its high transmission rates and its potential to cause severe, life-threatening complications in specific vulnerable populations, including neonates, pregnant women, adults, and immunocompromised individuals.

Although varicella is typically self-limiting in healthy children, the causative agent remains latent within the sensory nerve ganglia of the host for life. Years or decades later, the virus can reactivate, manifesting as herpes zoster (commonly referred to as shingles), which presents as a painful, localized dermatomal rash. Consequently, the clinical significance of varicella extends far beyond the acute childhood illness, encompassing long-term neuropathic risks and a substantial economic burden related to healthcare utilisation and lost productivity. In the modern medical era, the epidemiology and clinical impact of varicella have been database-proven and altered by the development and widespread implementation of highly effective live-attenuated vaccines, shifting the disease from an unavoidable childhood occurrence to a highly preventable vaccine-controlled condition.

\section*{Epidemiology}
Prior to the introduction of universal varicella vaccination programmes, varicella was an endemic disease with global distribution, affecting nearly every individual by early adulthood. In temperate climates, the disease exhibited a distinct seasonal pattern, peaking in late winter and early spring, whereas in tropical regions, transmission patterns were less seasonal, and individuals tended to acquire the infection at older ages. This age shift in tropical climates is hypothesized to be due to the heat-lability of the virus, which reduces its survival in high ambient temperatures, thereby delaying exposure and transmission to later childhood or adolescence.

The infectivity of varicella is exceptionally high, with secondary attack rates in susceptible household contacts approaching 90\%. The primary reservoir for the virus is humans, and there is no known animal vector. Transmission occurs primarily through the inhalation of airborne droplets or aerosols generated from the respiratory tract secretions of infected individuals, or via direct contact with the fluid of active vesicular skin lesions. An infected individual is contagious from approximately 48 hours before the onset of the characteristic rash until all lesions have crusted over, a period that typically spans 5--7 days.

Following the licensing and widespread implementation of the varicella vaccine in the mid-1990s, particularly in countries such as the United States, some countries, and parts of Europe, the epidemiology of the disease has changed dramatically. Incidences of varicella have declined by up to 90\% to 95\% in countries with high vaccine coverage, and there has been a corresponding, marked reduction in varicella-related hospitalizations and deaths. However, in regions where vaccination is not routinely administered, varicella continues to cycle epidemically, primarily affecting young children between the ages of 1--9 years. Adult varicella, while less common in endemic settings, represents a major clinical concern because the risk of severe complications, hospitalization, and mortality is 10- to 20-fold higher in adults than in children.

\section*{Aetiology and Pathophysiology}
The pathogenesis of varicella is a complex, multi-stage process initiated by exposure to the varicella-zoster virus. Below is a detailed breakdown of the causative agents, transmission routes, risk factors, and physiological pathways of the infection:

\begin{itemize}
    \item \textbf{Causative Agent}: The varicella-zoster virus (VZV) is an enveloped, double-stranded DNA virus belonging to the genus Varicellovirus of the subfamily Alphaherpesvirinae. The viral genome encodes approximately 70 unique proteins, many of which are involved in immune evasion, viral replication, and the establishment of latency.
    \item \textbf{Portal of Entry and Initial Replication}: The virus enters the host through the upper respiratory tract mucosa or the conjunctiva. Initial replication occurs in the regional lymphoid tissue of the nasopharynx, particularly the tonsils, over a period of 2--4 days.
    \item \textbf{Primary Viraemia}: Following initial local replication, the virus enters the bloodstream, resulting in a primary, transient viraemia. This occurs approximately 4--6 days after exposure and facilitates the transport of the virus to the reticuloendothelial system, including the liver, spleen, and sensory ganglia.
    \item \textbf{Secondary Viraemia}: A second wave of viral replication takes place within the internal organs over the next several days. This culminates in a more intense, secondary viraemia occurring 10--14 days post-exposure. During this phase, VZV is transported via infected T-lymphocytes (specifically skin-homing CD4+ and CD8+ memory T-cells) to the cutaneous capillary endothelial cells, leading to infection of the epidermis and the development of the characteristic skin lesions.
    \item \textbf{Establishment of Latency}: During the acute phase of varicella, the virus travels retrogradely along the sensory nerve fibers from the cutaneous lesions to the dorsal root, cranial nerve, or autonomic ganglia. VZV establishes a lifelong, latent infection within these neuronal cells, remaining clinically silent under the control of host cell-mediated immunity.
    \item \textbf{Risk Factors for Severe Disease}: While any non-immune individual can contract varicella, certain populations are at a significantly higher risk for severe pathophysiological progression and complications. These include \textbf{immunocompromised patients} (such as those with HIV/AIDS, malignancies, or those undergoing immunosuppressive therapy), \textbf{pregnant women} (due to altered cellular immunity and risks to the fetus), \textbf{neonates} (who have immature immune systems), and \textbf{chronic disease sufferers} (particularly those with underlying pulmonary or cutaneous conditions).
\end{itemize}

\section*{Clinical Features}
The clinical course of varicella typically begins after an incubation period of 14--16 days, though it can range from 10--21 days. The presentation varies in severity based on the age and immunological status of the patient. The clinical features manifest in distinct stages:

\begin{itemize}
    \item \textbf{Prodromal Phase}: In children, the prodrome may be absent or brief, lasting less than 24 hours. In adolescents and adults, a more prominent prodrome of 1--2 days is common, characterized by low-grade fever, malaise, myalgia, headache, loss of appetite, and mild abdominal pain.
    \item \textbf{Eruptive Phase (The Rash)}: The hallmark of varicella is the characteristic pruritic rash, which typically begins on the trunk and face before spreading centripetally to the scalp and extremities. The rash progresses rapidly through successive, distinct stages:
    \begin{itemize}
        \item \textbf{Macules}: The lesions begin as small, red, flat spots.
        \item \textbf{Papules}: Within hours, the macules evolve into raised, solid bumps.
        \item \textbf{Vesicles}: The papules quickly develop into thin-walled, fluid-filled blisters situated on an erythematous base, often described as resembling ``dewdrops on a rose petal.'' The fluid is initially clear but gradually becomes cloudy.
        \item \textbf{Pustules and Crusts}: The vesicles rupture, pustulate, and then dry out to form crusts or scabs within 24--48 hours. The crusts typically fall off within 1--2 weeks, leaving temporary hypopigmentation or hyperpigmentation, but no permanent scarring unless secondary bacterial infection occurs.
    \end{itemize}
    \item \textbf{Polymorphous Presentation}: A defining clinical feature of varicella is the presence of lesions in all stages of development simultaneously in any given area of the body. This occurs because the rash erupts in successive crops over a period of 3--5 days.
    \item \textbf{Mucosal Involvement}: Vesicular lesions frequently develop on mucosal surfaces, including the oral cavity, conjunctivae, pharynx, and vulvovaginal mucosa. These lesions rapidly ulcerate, causing significant discomfort, pain, and difficulty swallowing.
    \item \textbf{Systemic Symptoms}: Throughout the eruptive phase, patients experience mild to moderate fever (typically ranging from 37.8 degrees Celsius to 38.9 degrees Celsius, though it can be higher in severe cases) and generalized malaise. The severity of the fever usually correlates with the extent of the skin eruption.
\end{itemize}

\section*{Differential Diagnosis}
When evaluating a patient with a generalized vesicular rash, it is essential to consider a range of infectious and non-infectious conditions. The differential diagnosis for varicella includes the following:

\begin{itemize}
    \item \textbf{Disseminated Herpes Simplex Virus (HSV) Infection}: Can present with generalized vesicular lesions, particularly in patients with pre-existing skin conditions such as atopic dermatitis (eczema herpeticum) or in immunocompromised individuals. Unlike varicella, HSV lesions often cluster and are associated with different localized symptoms.
    \item \textbf{Hand, Foot, and Mouth Disease}: Caused by enteroviruses, most commonly Coxsackievirus A16 or Enterovirus A71. The vesicular rash is typically localized to the hands, feet, buttocks, and oral mucosa, and the lesions are usually less pruritic and more oval-shaped than those of varicella.
    \item \textbf{Impetigo}: A superficial bacterial skin infection, usually caused by Streptococcus pyogenes or Staphylococcus aureus. Bullous impetigo can present with fluid-filled blisters, but they are typically localized, lack the widespread centripetal distribution of varicella, and form characteristic honey-colored crusts.
    \item \textbf{Insect Bites}: Bites from fleas, bedbugs, or mites can cause pruritic papules and vesicles (papular urticaria). These are distinguished by their distribution (often clustered in linear patterns on exposed skin) and the absence of systemic prodromal symptoms or mucosal involvement.
    \item \textbf{Rickettsialpox}: A rare zoonotic infection caused by Rickettsia akari and transmitted by house mouse mites. It presents with a papulovesicular rash resembling varicella, but it is preceded by a diagnostic black eschar at the site of the mite bite.
    \item \textbf{Drug Eruptions}: Certain adverse drug reactions, such as generalized bullous fixed drug eruptions or early stages of Stevens-Johnson syndrome, can present with vesicular or bullous skin lesions. A detailed medication history and the presence of mucosal involvement without infectious prodrome are critical for differentiation.
    \item \textbf{Contact Dermatitis}: Acute allergic contact dermatitis (e.g., from poison ivy or chemical exposure) can produce highly pruritic, vesicular eruptions. These are typically localized to the area of contact and present in a linear or geometric pattern rather than a generalized, centripetal distribution.
    \item \textbf{Molluscum Contagiosum}: A viral skin infection caused by a poxvirus, producing small, firm, dome-shaped papules with central umbilication. Although they can persist for months, they lack the rapid evolution, fluid-filled vesicular nature, and systemic symptoms of acute varicella.
\end{itemize}

\section*{When to Seek Medical Care}
While varicella is often a self-limiting illness in healthy children, prompt medical evaluation is necessary in many scenarios to prevent severe complications. Patients, parents, and caregivers should seek medical attention under the following circumstances:

\begin{itemize}
    \item \textbf{High-Risk Individuals}: Any individual at high risk for severe disease who has been exposed to varicella or develops symptoms must contact a healthcare provider immediately. This includes pregnant women, newborns, infants under 1 year of age, adults, and individuals with compromised immune systems (such as those with cancer, HIV, or taking immunosuppressive drugs).
    \item \textbf{Signs of Secondary Bacterial Infection}: If the skin lesions become intensely red, swollen, warm, tender, or begin leaking yellow pus, it may indicate a secondary bacterial skin infection.
    \item \textbf{Neurological Changes}: Immediate medical evaluation is required if the patient exhibits extreme lethargy, confusion, difficulty walking, an unsteady gait (ataxia), a stiff neck, persistent vomiting, or seizures.
    \item \textbf{Respiratory Complications}: The development of a persistent cough, rapid breathing, wheezing, chest pain, or difficulty breathing (dyspnea) warrants urgent assessment, as these are signs of varicella pneumonia.
    \item \textbf{Severe Systemic Symptoms}: Seek care if the fever lasts longer than 4 days, rises above 39 degrees Celsius, or if the patient is unable to keep fluids down and shows signs of dehydration (e.g., decreased urination, dry mouth, extreme sleepiness).
\end{itemize}

For life-threatening symptoms, such as severe respiratory distress, unresponsiveness, or prolonged seizures, emergency services must be contacted immediately. In many countries, call local emergency services; in the United States, call 911; in the United Kingdom, call 999; or contact the relevant local emergency number in your jurisdiction.

\section*{Diagnosis and Investigations}
In the vast majority of cases, the diagnosis of varicella is clinical, based on the characteristic appearance and evolution of the skin rash combined with a history of exposure and prodromal symptoms. However, laboratory investigations are indicated in atypical presentations, hospitalized patients, pregnant women, neonates, and to confirm cases in highly vaccinated populations.

\begin{itemize}
    \item \textbf{Clinical Examination}: A thorough physical examination is performed to assess the morphology, stage, and distribution of the skin lesions, as well as to check for involvement of the mucous membranes and signs of systemic complications (such as auscultation of the lungs for signs of pneumonia).
    \item \textbf{Polymerase Chain Reaction (PCR)}: PCR testing is the gold standard and the most sensitive and specific laboratory method for confirming VZV infection. Swabs are taken from the base of an active vesicle, or crusts/scabs are collected and sent for analysis. PCR can rapidly detect VZV DNA and differentiate between wild-type virus and vaccine strains.
    \item \textbf{Direct Fluorescent Antibody (DFA) Staining}: This test involves collecting cells from the base of a vesicle and staining them with fluorescein-conjugated monoclonal antibodies directed against VZV glycoproteins. It provides rapid results but is less sensitive than PCR and requires an adequate cellular specimen.
    \item \textbf{Viral Culture}: VZV can be isolated in cell culture from vesicular fluid, but this method is slow (taking up to 1--2 weeks), technically demanding, and significantly less sensitive than PCR, as the virus is highly labile and difficult to culture.
    \item \textbf{Serological Testing}: Enzyme-Linked Immunosorbent Assay (ELISA) is commonly used to detect VZV-specific antibodies.
    \begin{itemize}
        \item \textbf{IgM Antibodies}: The presence of VZV-specific IgM antibodies indicates acute or very recent infection. However, IgM tests can sometimes yield false-positive or false-negative results.
        \item \textbf{IgG Antibodies}: A fourfold or greater rise in VZV-specific IgG antibody titres between acute and convalescent phase sera (collected 10--14 days apart) confirms acute infection. IgG testing is also widely used to assess immunity to VZV in high-risk individuals prior to exposure.
    \end{itemize}
    \item \textbf{Tzanck Smear}: A historical cytological test where scrapings from the base of a vesicle are stained (e.g., with Giemsa stain) to look for multinucleated giant cells and intranuclear inclusion bodies. While quick, it is non-specific as it cannot differentiate between VZV and herpes simplex virus infections.
\end{itemize}

\section*{Management}
The management of varicella is tailored to the age of the patient, their risk of complications, and the severity of their clinical presentation. It encompasses supportive measures, antiviral pharmacotherapy, and post-exposure prophylaxis.

\begin{itemize}
    \item \textbf{Supportive and Symptomatic Care}: For healthy children with mild disease, management is primarily supportive and aimed at relieving symptoms:
    \begin{itemize}
        \item \textbf{Fever Management}: Paracetamol (acetaminophen) is the preferred antipyretic for reducing fever and discomfort. \textbf{Aspirin and other salicylates must be strictly avoided} in children and adolescents with varicella due to the high risk of Reye's syndrome, a rare but life-threatening condition characterized by acute encephalopathy and fatty liver infiltration. Non-steroidal anti-inflammatory drugs (NSAIDs), such as ibuprofen, are also generally avoided by many clinicians due to epidemiological associations with an increased risk of severe secondary bacterial skin infections (necrotizing fasciitis).
        \item \textbf{Pruritus Relief}: Cool baths, often with added colloidal oatmeal, can soothe itchy skin. Topical calamine lotion may be applied to individual lesions. Oral antihistamines (such as cetirizine or diphenhydramine) can help reduce itching and facilitate sleep.
        \item \textbf{Hydration and Nutrition}: Maintaining adequate fluid intake is essential, particularly if oral mucosal lesions make eating and drinking painful. Soft, bland foods are recommended to avoid irritating oral ulcers.
    \end{itemize}
    \item \textbf{Antiviral Therapy}: Antiviral medications (most commonly acyclovir or valacyclovir) inhibit viral replication and are indicated for individuals at risk of moderate to severe disease:
    \begin{itemize}
        \item \textbf{Oral Antivirals}: Recommended for healthy adolescents and adults, patients with chronic cutaneous or pulmonary diseases, individuals receiving long-term salicylate therapy, and those taking short courses of corticosteroids. To be effective, oral antiviral therapy must be initiated within 24 hours of rash onset.
        \item \textbf{Intravenous Antiviral Therapy}: Reserved for patients with severe disease, complications (such as varicella pneumonia or encephalitis), and immunocompromised individuals. Intravenous acyclovir is administered in a hospital setting to ensure therapeutic blood levels and monitor renal function.
    \end{itemize}
    \item \textbf{Post-Exposure Prophylaxis}: For susceptible individuals at high risk of severe varicella who have had close exposure to the virus:
    \begin{itemize}
        \item \textbf{Varicella-Zoster Immunoglobulin (VZIG)}: A passive immunization product prepared from human plasma with high titres of VZV antibodies. It should be administered intramuscularly as soon as possible, and ideally within 96 hours (up to 10 days) of exposure, to prevent or attenuate the infection.
        \item \textbf{Post-Exposure Vaccination}: For healthy, unvaccinated individuals without evidence of immunity, administering the varicella vaccine within 3--5 days of exposure can prevent infection or modify the severity of the disease.
    \end{itemize}
\end{itemize}

\section*{Complications}
While varicella is often benign, it can lead to serious, life-threatening complications. The risk is elevated in infants, adults, pregnant women, and the immunocompromised. The primary complications include:

\begin{itemize}
    \item \textbf{Secondary Bacterial Infections}: The most common complication of varicella in children. Scratching the pruritic lesions compromises the skin barrier, allowing bacteria to enter. The primary pathogens are \textit{Staphylococcus aureus} and \textit{Streptococcus pyogenes} (Group A Streptococcus). These infections can range from mild superficial impetigo or cellulitis to severe, invasive infections such as necrotizing fasciitis, osteomyelitis, septicaemia, and streptococcal toxic shock syndrome.
    \item \textbf{Neurological Complications}: VZV has neurotropic properties, and nervous system involvement can manifest in several ways:
    \begin{itemize}
        \item \textbf{Acute Cerebellar Ataxia}: A relatively common, benign complication in children, typically occurring 1--3 weeks after the rash. It presents with an unsteady gait, nystagmus, and slurred speech, and usually resolves spontaneously within weeks to months.
        \item \textbf{Encephalitis}: A rare but severe complication, more common in adults, presenting with confusion, fever, seizures, and altered consciousness. It requires intensive care and intravenous antiviral therapy.
        \item \textbf{Other Neurological Manifestations}: These include transverse myelitis, Guillain-Barré syndrome, aseptic meningitis, and cranial nerve palsies.
    \end{itemize}
    \item \textbf{Varicella Pneumonia}: The most common serious complication of varicella in adults and pregnant women, accounting for significant morbidity and mortality. It typically develops 1--6 days after the rash onset and is characterized by cough, dyspnea, tachypnea, chest pain, and hemoptysis. Radiography typically reveals diffuse, nodular infiltrates.
    \item \textbf{Congenital Varicella Syndrome}: Occurs when a pregnant woman contracts varicella during the first 20 weeks of gestation. The virus crosses the placenta, leading to fetal infection that can cause intrauterine growth restriction, cicatricial skin scarring, limb hypoplasia, microcephaly, chorioretinitis, cataracts, and severe developmental delays.
    \item \textbf{Neonatal Varicella}: A highly severe form of the disease occurring in newborns whose mothers develop varicella rash within 5 days before to 2 days after delivery. Because the mother has not had time to develop and pass protective IgG antibodies to the fetus, the newborn is exposed to high viral loads without passive immunity, resulting in an attack rate of up to 30\% and a high mortality rate if untreated.
    \item \textbf{Other Complications}: These include thrombocytopenia (which can lead to epistaxis, petechiae, or purpura fulminans), glomerulonephritis, myocarditis, pancreatitis, orchitis, and transient arthritis.
\end{itemize}

\section*{Prognosis and Outcomes}
For the vast majority of healthy children, the prognosis of varicella is excellent. The illness is typically self-limiting, with acute symptoms resolving within 7--10 days. The cutaneous lesions heal completely without scarring, provided there is no secondary bacterial infection or severe mechanical trauma from scratching.

However, the prognosis is significantly more guarded in high-risk populations. In healthy adults, the risk of hospitalization is approximately three times higher than in children, and the mortality rate is up to 20 times higher. For pregnant women, varicella carries not only the risk of maternal complications, such as life-threatening varicella pneumonia, but also the risk of congenital varicella syndrome or neonatal varicella, which carry substantial long-term morbidity and mortality. In immunocompromised individuals, varicella can progress to a severe, disseminated disease characterized by visceral involvement, haemorrhagic rash, and multi-organ failure, with mortality rates historically reaching 15\%--18\% in the absence of prompt antiviral therapy.

A critical aspect of the long-term prognosis of varicella is the persistence of VZV in a latent state within the sensory ganglia. Approximately one in three individuals who have had varicella will experience reactivation of the virus later in life, resulting in herpes zoster (shingles). Herpes zoster is characterized by a painful, localized vesicular rash, and its most common complication is postherpetic neuralgia (PHN), a debilitating chronic neuropathic pain syndrome that can persist for months or years.

\section*{Prevention and Self-Care}
The primary strategy for controlling varicella is active immunization, complemented by appropriate home care and infection control measures.

\begin{itemize}
    \item \textbf{Active Immunization}:
    \begin{itemize}
        \item \textbf{Varicella Vaccine}: A live-attenuated vaccine (e.g., Oka strain) is highly effective at preventing varicella. It is typically administered as a two-dose series. In many childhood immunization schedules, the first dose is given at 12--15 months of age, and the second dose at 4--6 years of age. A combination vaccine containing measles, mumps, rubella, and varicella (MMRV) is also commonly used.
        \item \textbf{Adolescent and Adult Vaccination}: Susceptible individuals aged 13 years and older who have no evidence of immunity should receive two doses of the varicella vaccine, administered 4--8 weeks apart.
        \item \textbf{Herd Immunity}: High vaccine coverage in the community protects vulnerable individuals who cannot receive the live-attenuated vaccine (e.g., pregnant women and severely immunocompromised individuals) by reducing the overall circulation of the virus.
    \end{itemize}
    \item \textbf{Infection Control}:
    \begin{itemize}
        \item \textbf{Isolation}: To prevent transmission, individuals with active varicella must be isolated from school, work, and public settings until all vesicles have completely crusted over.
        \item \textbf{Avoid Contact with Vulnerable Groups}: Special care must be taken to prevent contact between infected individuals and susceptible pregnant women, neonates, or immunocompromised individuals.
    \end{itemize}
    \item \textbf{Home Care and Self-Care Measures}:
    \begin{itemize}
        \item \textbf{Preventing Scratching}: Keep fingernails clean and trimmed short. Mittens or socks can be worn on the hands of young children to prevent scratching, which reduces the risk of secondary bacterial infections and scarring.
        \item \textbf{Skin Care}: Bathe daily in lukewarm water. Adding colloidal oatmeal or baking soda to the bath can help relieve pruritus. Pat the skin dry with a clean towel rather than rubbing.
        \item \textbf{Clothing}: Wear loose-fitting, soft cotton clothing to minimize irritation to the skin lesions.
        \item \textbf{Hydration}: Encourage frequent intake of cool fluids (water, diluted juices, ice pops) to prevent dehydration.
        \item \textbf{Oral Hygiene}: If oral lesions are present, offer a soft diet and avoid salty, acidic, or spicy foods. Rinsing the mouth with salt water (half a teaspoon of salt in a glass of warm water) can soothe discomfort in older children and adults.
    \end{itemize}
\end{itemize}

\section*{Sources and Further Reading}
For further information, guidelines, and updates on varicella and vaccination protocols, please refer to the following authoritative resources:

\begin{itemize}
    \item World Health Organization (WHO): \texttt{https://www.who.int}
    \item Centers for Disease Control and Prevention (CDC): \texttt{https://www.cdc.gov}
    \item Australian Government Department of Health and Aged Care: \texttt{https://www.health.gov.au}
    \item National Health Service (NHS): \texttt{https://www.nhs.uk}
    \item Johns Hopkins Medicine: \texttt{https://www.hopkinsmedicine.org}
    \item Mayo Clinic: \texttt{https://www.mayoclinic.org}
\end{itemize}